% Follow-up to H9-T6, independently checked by /root/programme_audit.
% Run JAC1964-20261121; GPT-6 Astra Pro; effort setting not exposed.

\subsection*{H9-T6a. A concrete generic-pencil polarization theorem}

\begin{theorem}[H9-T6a]
Let $X$ be smooth, projective and connected of even dimension $2m$
over an algebraically closed field $k$ of arbitrary characteristic.
Assume that the Lefschetz lowering operator of $X$ is algebraic, and
let $Y_\eta$ be the geometric generic fibre of a Lefschetz pencil of
hyperplane sections for a projective embedding of $X$.
The variable incidence factor of $P^m(Y_\eta)$ is polarized by its
$(-1)^m$-signed incidence correspondence.

If $X$ is abelian, then, accepting 032 Corollary~8.4, the same statement
holds for the entire primitive factor of $P^m(Y_\eta)$.
If $X$ is a smooth complete intersection in projective space, it holds
for all of $P^m(Y_\eta)$, without requiring 032.

For any starting projective embedding, a general pencil of degree-$d$
hypersurface sections is a Lefschetz pencil for every $d\ge2$.
Consequently these assertions give concrete families in every
characteristic, including characteristic two.
\end{theorem}

\begin{proof}
Write $n=\dim X=2m$ and $h$ for the hyperplane class of the pencil.
Hard Lefschetz gives a nondegenerate ambient summand
$j^*H^{n-1}(X)$ in $H^{n-1}(Y_\eta)$.  Its orthogonal complement is
the variable space $E_\ell=\ker j_*$.  The algebraic correspondence
\[
 1-j^*(L_h:H^{n-1}(X)\longrightarrow H^{n+1}(X))^{-1}j_*
\]
is the self-adjoint projector onto $E_\ell$ on the relevant degree.
Lemma~\ref{jac:P01} and Saito's transpose compatibility give a
nondegenerate incidence factor $B$ and a normalized rational
Poincar\'e cycle.  If $B=0$, its polarization assertion is empty.
Suppose $B\ne0$.

The incidence formula makes the realization map
$V_\ell B\to E_\ell(m)$ injective for every invertible prime $\ell$,
so $E_\ell\ne0$ for every such prime.  SGA~7~II, XVIII,
Theorem~6.6, printed pp.~323--325, says that the local vanishing
cycles span $E_\ell$ and are conjugate up to sign.  Thus every local
vanishing cycle is nonzero, for every $\ell$.  In particular the
critical set is nonempty and one can choose the same critical point
for all primes.  Corollary~6.7.2, printed p.~326, says that every
finite-index subgroup of geometric monodromy acts absolutely
irreducibly on $E_\ell$.  Its hypotheses are satisfied: $X$ is smooth
projective, the pencil is Lefschetz, hard Lefschetz holds for $X$, and
$n$ is even.  The exception in Theorem~6.6 is the conjunction
``$n$ odd and characteristic two'', which is absent here.

These are exactly the hypotheses used in H9-T6.  Its all-primes
monodromy argument therefore proves that the signed variable class
is a rational polarization.  A separate condition~A assumption was
unnecessary: nonzero $B$ itself supplied the nonzero variable
cohomology and hence the nonzero local vanishing cycles.

If $X$ is abelian, H9-T4 polarizes the fixed primitive summand by 032
on $J(C)\times X$.  The primitive fixed and variable projectors are
orthogonal, so their polarizations combine and descend under the
isogeny decomposition.  If $X$ is a smooth complete intersection of
even dimension, $H^{n-1}(X)=0$, so there is no fixed odd summand.
Moreover $H^{n+1}(Y_\eta)=0$ by dual weak Lefschetz, so all of its
middle odd cohomology is primitive.  Algebraicity of the lowering
operator for a smooth complete intersection is supplied explicitly
in JAC-P08 below.  Thus the variable conclusion is the entire
assertion in this case.

Finally, SGA~7~II, XVII, Theorem~2.5.1, printed p.~217, states that
every degree-$d$ multiple of a projective embedding, for $d\ge2$,
is a Lefschetz embedding.  The theorem has no characteristic
restriction.  The term ``quadratic ordinary'' includes the usual
characteristic-two distinction, but in our even ambient dimension
the local quadratic singularities are nondegenerate: XVII, p.~219
explicitly notes this when the ambient dimension is even or the
characteristic is not two.  This verifies that the odd-fibre
Picard--Lefschetz formula used in H9-T6 applies here as stated.
\end{proof}

\paragraph{The finite-index step can be checked without a density theorem.}
The proof of XVIII~6.7.2, printed p.~327, is short.  With the
alternating middle pairing, each local inertia generator acts by
$T_i=1+N_i$, where $N_i^2=0$ and
$\operatorname{im}N_i=\mathbf Q_\ell\delta_i$.
For any finite-index subgroup $G'$ a positive power $T_i^{a_i}$
belongs to its image, and
$T_i^{a_i}-1=a_iN_i$.  Since $a_i$ is nonzero in $\mathbf Q_\ell$,
a subspace stable under $G'$ is stable under every $T_i$.
The inertia groups topologically generate the tame monodromy group,
so it is stable under the full group.  Given a nonzero vector, the
nondegenerate pairing and the spanning vanishing cycles produce
some $N_i$ with nonzero image; conjugacy then gives every vanishing
cycle, hence the entire space.  The argument survives scalar
extension to $\overline{\mathbf Q}_\ell$, proving absolute
irreducibility.  No orthogonal-monodromy assertion is used here.

\begin{proposition}[JAC-P07: condition A is automatic for abelian pencils]
\label{jac:P07}
Let $X$ be an abelian variety of even dimension $n=2m\ge2$ over an
algebraically closed field, and let $D$ be a genuine Lefschetz pencil
for a projective embedding of $X$.  Let $L$ be its ample section
line bundle, $h=c_1(L)$, and $Y$ a smooth member.  For every prime
$\ell$ invertible in the field, the variable middle cohomology has
dimension
\[
 \dim E_\ell=\deg h^n-C_n,
 \qquad C_n=\frac1{n+1}\binom{2n}{n}.
                                                        \tag{JAC.P07.1}
\]
This dimension is strictly positive.  Every local vanishing cycle
is nonzero and the pencil satisfies condition~A of SGA~7~II,
XVIII~5.3, in every characteristic.  In particular the
specialization map $H^n(Y_s,\mathbf Q_\ell)\to
H^n(Y_{\overline\eta_s},\mathbf Q_\ell)$ at every singular fibre
is injective.
\end{proposition}

\begin{proof}
The tangent bundle of $X$ is trivial.  The Euler-characteristic
formula SGA~7~II, XVIII~6.4.2.3, printed p.~321, therefore gives
\[
 \chi(Y,\mathbf Q_\ell)=(-1)^{n-1}\deg h^n=-\deg h^n.
\]
Weak Lefschetz, Poincar\'e duality on $Y$, and
$b_i(X)=\binom{2n}{i}$ give
\[
 b_{n-1}(Y)=\deg h^n+
       2\sum_{i=0}^{n-2}(-1)^i\binom{2n}{i}
 =\deg h^n+2\binom{2n-1}{n-2}.
\]
The ambient subspace has dimension $b_{n-1}(X)=\binom{2n}{n-1}$.
Subtracting and using Pascal's identity yields (JAC.P07.1).
All identities are uniform in $\ell$.

Riemann--Roch on $X$ gives $\deg h^n=n!\chi(L)$, and for an
ample line bundle on an abelian variety $\chi(L)=h^0(L)$ is a
positive integer.  If $n\ge4$, then $n!>C_n$.  One may check this
without an estimate: $3!=6>C_3=5$, and
\[
 \frac{C_{r+1}}{C_r}=\frac{2(2r+1)}{r+2}<4\le r+1
 \qquad(r\ge3),
\]
so the inequality propagates.  If $n=2$, a genuine pencil requires
two independent sections of $L$, hence $h^0(L)\ge2$ and
$\deg h^2=2h^0(L)\ge4>C_2=2$.  This proves strict positivity
in every even dimension under consideration.

By XVIII~6.6.1--6.6.2, the local vanishing cycles span $E_\ell$
and are conjugate up to sign.  Since $E_\ell\ne0$, none of these
cycles is zero, and there is at least one critical fibre.
Proposition~6.2, printed p.~319, identifies this condition with
condition~A.  The excluded parity/characteristic combination in
XVIII~6.6 is odd ambient dimension in characteristic two; our
ambient dimension is even.  Thus the conclusion includes
characteristic two.

For clarity, the specialization consequence can also be seen from
the vanishing-cycle sequence at an ordinary quadratic fibre.
Its potentially nonzero connecting portion is
\[
 H^{n-1}(Y_{\overline\eta_s})\longrightarrow
 \mathbf Q_\ell(-m)\longrightarrow H^n(Y_s)
 \longrightarrow H^n(Y_{\overline\eta_s})\longrightarrow0.
\]
The first arrow is pairing with the nonzero local vanishing cycle,
up to the fixed nonzero normalization, and is surjective.
Consequently the final specialization map is injective (indeed an
isomorphism).  These are exactly the local hypotheses used in
H9-T2; integrality of the singular fibres alone would not suffice.
\end{proof}

\begin{proposition}[JAC-P08: further ambient varieties for Theorem~6]
\label{jac:P08}
The rational hyperplane and trace conclusions of JAC-P02 apply to
smooth complete intersections and to products of smooth complete
intersections, curves, and abelian varieties, in every characteristic.
For an even-dimensional smooth complete intersection $X$, the
critical incidence variety $P^m(Y_\eta)$ has zero $K/k$-trace up to
isogeny.  Products with nonzero ambient odd cohomology give further
instances of the nonzero fixed-part equality in JAC-P02.
\end{proposition}

\begin{proof}
For a smooth complete intersection $X$ of dimension $n$, with
hyperplane class $h$ and $e=\deg h^n>0$, weak Lefschetz and duality
show that its cohomology outside the middle degree consists of the
one-dimensional spaces generated by $h^r$ in degree $2r$.
The normalized lowering operator is the algebraic correspondence
\[
 \Lambda_X=\frac1e\sum_{r=1}^{n}
           h^{n-r}\times h^{r-1}
 \quad\in\mathrm{CH}^{n-1}(X\times X)_{\mathbf Q}.
\]
It sends $h^r$ to $h^{r-1}$.  It kills the middle primitive
cohomology: this is automatic by parity in odd dimension, and in
even dimension follows from its orthogonality to $h^{n/2}$.
The same elementary construction applies to curves.  For abelian
varieties algebraicity of the lowering operator is the known
Lefschetz standard theorem used throughout this paper.

For completeness, algebraic Lefschetz operators are stable under
products.  Replace each normalized lowering operator by the usual
$\mathfrak{sl}_2$ lowering operator, which acts on a primitive string
by $L^av\mapsto a(d-i-a+1)L^{a-1}v$.  The degree projectors and
primitive-string projectors are algebraic, so this replacement is
algebraic.  Explicitly, put $P=1-L\Lambda$. On primitive vectors the
operator
$P(\Lambda^rL^r-\Lambda^{r+1}L^{r+1})P$ selects precisely strings
of length $r+1$; conjugating these selectors by $L^a$ and
$\Lambda^a$ selects each level. This also constructs the needed
grading projectors directly from the two algebraic operators.
On a product the sums of the raising, lowering, and
grading operators give the tensor-product $\mathfrak{sl}_2$ action.
Its Casimir and grading have finite integer spectra.  Rational
polynomial interpolation in these commuting operators rescales the
nonzero lowering coefficients to one, and gives zero on primitive
vectors; hence the normalized product lowering operator is again
algebraic.  Applying JAC-P02 proves the assertions.

In the even complete-intersection case all odd cohomology of $X$
vanishes.  Lemma~\ref{jac:P01}, applied to the identity correspondence,
then gives $P^p(X)=0$ for every $p$.  The critical trace equality
of JAC-P02 therefore gives zero trace for $P^m(Y_\eta)$, even when
that incidence variety itself is nonzero.  For example, this is a
substantive assertion for generic curve sections of smooth surfaces
and for odd-dimensional generic complete intersections carrying a
nonzero middle incidence Picard variety.
\end{proof}

For a degree-$d$ hypersurface section relative to a starting class
$h$, formula (JAC.P07.1) reads
$\dim E_\ell=d^n\deg h^n-C_n$.  The dimension calculation is not
needed for the variable-polarization argument in H9-T6a, which can
start from $B\ne0$.  It is needed to discharge condition~A uniformly
in the geometric height theorem, including when $B=0$.
